% ═══════════════════════════════════════════════════════════════════════════════
\section{Literature Review: Re-evaluating GVCs, the State, and Structural Transformation}
\label{sec:literature_review}
% ═══════════════════════════════════════════════════════════════════════════════

The analytical framework of this dissertation draws upon three interconnected, and often competing, streams of political economy and development literature: (i) Global Value Chain (GVC) theory and the mechanics of industrial upgrading \cite{10_1038_s41592_019_0686_2, 10_1016_s0140_6736(20)30566_3, 10_1016_s0140_6736(20)30752_2, gereffi2005, milberg2013}; (ii) developmental state theory, the East Asian industrialisation experience, and the rebirth of industrial policy \cite{10_1038_s41586_020_2649_2, 10_1016_j_lindif_2023_102274, 10_1063_5_0007045, amsden1989, wade1990}; and (iii) the structural constraints facing landlocked developing countries (LLDCs) within the context of Eurasian integration and the Belt and Road Initiative (BRI) \cite{10_46328_ijonse_32, 10_1016_j_scs_2020_102382, 10_3390_w12102908, 10_1017_9781009157964_012, 10_1186_s13705_019_0232_1}.

By synthesising these theoretical domains, this chapter establishes the conceptual foundation for evaluating whether Uzbekistan's post-2017 liberalisation trajectory represents a genuine structural transformation or merely a reconfiguration of its historical commodity dependence \cite{10_1016_s0140_6736(21)02796_3, 10_1016_s0140_6736(24)00550_6}. It critically examines the contradictions inherent in pursuing state-led development within deeply integrated, transnational production regimes.

\subsection{Global Value Chains: Governance, Upgrading, and Contemporary Reconfigurations}

The globalisation of production networks over the past four decades has fundamentally altered the paradigm of industrial development \cite{10_1038_s41467_021_21246_9, 10_1186_s40537_021_00444_8, 10_1017_rdc_2020_41}. Rather than building entire, vertically integrated domestic supply chains from raw materials to finished goods, developing economies now pursue industrialisation by inserting themselves into specific nodes of global and regional production networks \cite{gereffi2005, milberg2013, 10_1136_bmj_n2061, 10_1038_s41586_020_2286_9}. 

A critical contribution of GVC theory is the concept of ``governance''---the power dynamics dictating how financial, material, and human resources are allocated within a chain. \citet{gereffi2005} distinguish between five governance structures: market, modular, relational, captive, and hierarchy. The upgrading trajectory of a developing country depends heavily on the governance structure it enters \cite{10_1007_s00134_021_06506_y, 10_1561_2200000083}. \citet{humphrey2002} demonstrate that captive governance---where powerful lead firms strictly control specifications and limit technology diffusion---often facilitates rapid product and process upgrading but actively hinders functional upgrading (moving into design, marketing, or independent branding) \cite{10_1111_brv_12627, 10_37074_jalt_2020_3_1_7}. Conversely, relational governance fosters deeper knowledge exchange \cite{10_1109_comst_2023_3249835, 10_1093_qje_qjz041}. 

The ``smile curve'' phenomenon \cite{baldwin2016} further complicates GVC insertion \cite{10_1109_access_2021_3140175, 10_5465_amr_2018_0072}. Value capture has increasingly concentrated at the pre-production (R\&D, design) and post-production (marketing, retail) ends of the curve, while the middle stages (manufacturing and assembly) are heavily commoditised. Consequently, \citet{rodrik2016} warns of ``premature deindustrialisation,'' suggesting that the classical pathway of absorbing surplus agricultural labour into low-skilled manufacturing is closing \cite{10_1109_access_2020_2970143, 10_1038_s41560_019_0513_0}. 

Recent literature has increasingly focused on the vulnerability and reconfiguration of these chains \cite{10_1007_s11023_020_09517_8, 10_3145_epi_2020_ene_03}. \citet{10_1057_s42214_020_00062_w} highlights how recent global shocks have exposed the fragility of hyper-optimised GVCs, prompting lead firms to prioritise resilience over cost-efficiency \cite{10_1016_j_ijinfomgt_2021_102456, 10_1016_j_socscimed_2020_113370}. This has accelerated ``China+1'' strategies and regionalisation \cite{10_1057_s42214_021_00102_z, 10_1007_s10842_019_00322_3, 10_1080_0960085x_2020_1718007}. For Uzbekistan, the theoretical imperative is to leverage this regionalisation to secure foreign direct investment (FDI) that facilitates progression from captive assembly to relational upgrading \cite{gereffi2016, 10_3133_mcs2021, 10_1162_coli_a_00502}, particularly by embedding itself into the shifting Northeast Asian production networks. Furthermore, \citet{10_1057_s42214_020_00071_9} emphasise that understanding regional value chains requires analysing the interaction between private lead-firm governance and public institutional frameworks \cite{10_1016_j_resconrec_2021_105618, 10_1038_s41893_021_00784_6}.

\subsection{The Developmental State and the Rebirth of Industrial Policy}

The integration of Uzbekistan into Northeast Asian networks must be analysed through the theoretical lens of East Asian industrialisation \cite{10_1109_tro_2021_3075644, 10_1016_s0140_6736(21)01330_1}. The classical interpretation of the East Asian ``miracle'' rests on the concept of the developmental state, which challenges the neoclassical orthodoxy that East Asian growth resulted primarily from free markets and minimal state intervention \cite{amsden1989, wade1990, evans1995, 10_1016_j_neucom_2020_07_061}.

\citet{amsden1989} demonstrated how South Korea's industrialisation relied on ``getting prices wrong''---using targeted subsidies, directed credit, and strategic protectionism to incubate domestic champions. \citet{wade1990} advanced the ``governed market'' theory, illustrating how Taiwan's state apparatus actively steered investment into strategic sectors \cite{10_1007_s11157_020_09523_3, 10_18260_1_2__22585}. \citet{evans1995} introduced the concept of ``embedded autonomy,'' arguing that successful developmental states require bureaucracies that are sufficiently insulated from rent-seeking pressures (autonomy) yet deeply connected to the industrial capitalist class to facilitate policy implementation (embeddedness) \cite{10_1038_s41746_020_00323_1, 10_1038_s41558_020_0797_x}. 

\citet{chang2002} famously argued that developed nations often ``kick away the ladder'' by denying developing countries the same protectionist tools they used to industrialise \cite{10_1111_dech_12593}. However, recent literature indicates a ``rebirth'' of industrial policy globally \cite{10_1016_j_respol_2023_104863}. \citet{10_1007_s10842_019_00322_3} argues for a new, systemic industrial policy focused on societal goals rather than merely sectoral subsidies. \citet{10_1080_01436597_2023_2217766} observes the advance of the state in the age of strategic competition, where industrial policy is renewed as a tool for geopolitical positioning as much as economic development \cite{10_1088_1748_9326_abc9e1}. \citet{lin2012} attempts to reconcile state intervention with market forces through ``New Structural Economics,'' advocating for state facilitation that aligns with a country's latent comparative advantage \cite{10_1590_0101_31572020_3102, 10_1057_s41287_020_00256_1}. 

For Uzbekistan, the challenge is constructing an institutional apparatus capable of exercising embedded autonomy within the constraints of modern trade rules, navigating the tension between required liberalisation and the necessary state intervention to foster domestic capabilities \cite{10_1093_ia_iiaa064, 10_3389_fenvs_2022_938776, 10_1038_s41467_020_17710_7}. Institutional readiness remains a bottleneck \cite{10_1111_gcb_15873, 10_1186_s13045_021_01213_z, 10_3133_mcs2020}.

\subsection{The Flying Geese Paradigm and South-South Cooperation}

The spatial dynamics of East Asian FDI are elegantly captured by the ``flying geese'' model \cite{akamatsu1962, 10_3390_su12051859}. The model describes a tiered hierarchy in which a lead economy transfers sunset, labour-intensive industries to follower economies as its own comparative advantage shifts toward capital- and knowledge-intensive production \cite{dunning1981, 10_1057_s41599_023_02234_4, 10_1590_0034_7329202000201}. 

Applying the flying geese model to contemporary Central Asia requires significant modification. While China is graduating from labour-intensive manufacturing, creating a theoretical opportunity for Uzbekistan \cite{10_1016_j_jclepro_2020_122263}, Chinese outward FDI often diverges from the classical model. Scholars analyse the structural transformation of Chinese South-South cooperation, noting that it often prioritises infrastructure and resource extraction over the export-platform manufacturing seen in the traditional flying geese cascade \cite{10_1080_00213624_2022_2020025, 10_1007_s12571_021_01184_6, 10_3390_su13020928, 10_1038_s41467_022_29324_2}. The Belt and Road Initiative (BRI), while providing critical infrastructure \cite{10_1017_s0305741023000966, 10_1057_s41278_020_00180_5}, may deepen primary commodity dependence if not paired with robust domestic industrial policies \cite{10_1002_jid_3664, 10_1038_s41893_022_00909_5}.

\subsection{Comparative Development Benchmarks: Vietnam and Thailand}

Uzbekistan's aspirations can be benchmarked against two successful models of GVC insertion: Thailand and Vietnam \cite{10_1016_j_techfore_2020_119937, 10_3390_land9080275}. Thailand's Eastern Seaboard Development Programme (1980s--1990s) leveraged Japanese FDI following the 1985 Plaza Accord \cite{10_1002_joc_6607, 10_1111_1468_2427_13080}. The state successfully utilised local content requirements to force technology transfer, establishing Thailand as a regional automotive hub \cite{10_1007_s11069_021_05150_5}. 

Vietnam represents a contemporary model of rapid GVC embedding following the Doi Moi reforms and its 2007 WTO accession \cite{10_1016_j_ecoenv_2023_114665, 10_1007_s00376_025_4540_4, 10_1038_s42949_020_00012_8}. Vietnam successfully positioned itself as a ``China+1'' destination, attracting massive electronics assembly FDI \cite{10_3390_land11122253, 10_1371_journal_pone_0243097}. Vietnam's trajectory underscores the importance of deep trade liberalisation, proactive investment promotion, and SEZ policy \cite{10_1136_bjsports_2020_102955, 10_1016_s1474_4422(21)00252_0, 10_1016_s0140_6736(20)30045_3}.

\subsection{Central Asian Structural Constraints and Landlocked Theory}

Uzbekistan's integration prospects are uniquely bounded by its physical geography \cite{10_3389_fsoc_2021_629693, 10_5089_9781513595405_001}. \citet{limao2001} demonstrate that landlocked countries face a significant structural penalty, with transport costs averaging 50\% higher than coastal economies \cite{10_1016_j_apenergy_2020_115026, 10_1177_11786221211028185}. For a doubly landlocked nation, this geographical penalty precludes integration into time-sensitive, low-margin GVCs \cite{10_1038_s41467_022_32070_0, 10_3389_fmars_2020_00586}.

These structural constraints dictate a specific industrialisation pathway. High-volume, low-value assembly is economically unviable due to prohibitive transit costs \cite{10_5089_9798400249006_001, 10_1016_j_apenergy_2021_117948}. Industrial upgrading must target goods with high value-to-weight ratios or leverage regional markets where transit penalties are minimised. The political economy of the Middle Corridor and the China-Kyrgyzstan-Uzbekistan railway are therefore existential prerequisites for Uzbekistan's GVC embedding \cite{10_1007_s10666_021_09750_0, 10_1038_s43016_022_00588_7, 10_1016_j_egyr_2022_01_200, 10_1038_s41558_021_01219_y}. Recent literature on logistics and Eurasian corridors highlights that multimodal transport efficiency is the new driver of competitiveness \cite{10_1038_s41586_022_05453_y, 10_1057_s41267_021_00435_0, 10_1007_s10708_020_10320_2}.

\subsection{Technology Transfer and FDI Spillovers}

The final theoretical piece concerns how FDI translates into domestic capability \cite{10_1257_aer_20180338, 10_1007_978_3_030_90434_0_43_1}. Theoretical models of FDI spillovers suggest that foreign presence can increase domestic productivity through demonstration effects, labour mobility, and vertical linkages (both forward and backward) \cite{10_3390_su132111663, 10_1007_s11356_021_15304_4}. However, empirical evidence is mixed. \citet{10_1007_s11356_020_08812_2} show that without sufficient ``absorptive capacity'' (human capital, R\&D investment) in domestic firms, foreign presence can actually crowd out local enterprises \cite{10_1017_s0140525x22002023}. Thus, SEZ policies in Uzbekistan must be explicitly tied to domestic supply-chain integration rather than operating as isolated enclaves \cite{10_21832_9781800417151_006, 10_15585_mmwr_ss6904a1, 10_1016_s0140_6736(24)01867_1}.

\subsection{Synthesis: Economic Complexity and Upgrading}

To measure whether these varied policy efforts are translating into genuine structural transformation, this study relies on the Economic Complexity Index (ECI) \cite{hidalgo2009, hausmann2014}. The ECI posits that development is fundamentally a process of accumulating productive knowledge \cite{10_1093_jeea_jvab012}. \citet{10_1016_j_respol_2023_104863} further elaborates on the policy implications of economic complexity, noting that countries must strategically target related, more complex industries \cite{10_1007_s11187_021_00544_y, 10_1016_j_enpol_2021_112322}. The empirical analysis that follows will assess whether Uzbekistan's integration into Northeast Asian networks is driving capability accumulation, moving the nation beyond its historical dependency on resource extraction \cite{10_1016_j_eneco_2020_105071, 10_1111_anti_12699, 10_1057_s41278_024_00287_z}.

